\section{Related Work}
\label{sec:benchmarks_related_work}

This section reviews prior work on evaluation environments and benchmarks relevant to the study of agents under open-world novelty. The discussion is organized around four themes: standard reinforcement learning and planning benchmarks that operate under closed-world assumptions, environments that introduce variability or non-stationarity, frameworks specifically designed for open-world novelty evaluation, and evaluation protocols and metrics for measuring agent adaptation.

\paragraph{Closed-world evaluation benchmarks.}
The dominant paradigm for evaluating sequential decision-making agents assumes fixed, fully specified environments. Reinforcement learning benchmarks such as OpenAI Gym~\citep{brockman2016openai} and the Arcade Learning Environment~\citep{bellemare2013arcade} provide standardized tasks with known dynamics, enabling comparison of learning algorithms under controlled conditions. Offline RL benchmarks such as D4RL~\citep{fu2020d4rl} extend this paradigm to data-driven evaluation but retain the assumption of a stationary environment. Similarly, classical planning benchmarks expressed in PDDL~\citep{mcdermott1998pddl} assume that the domain model is complete and correct, and evaluate planners on plan quality and computational efficiency. PDDLGym~\citep{silver2020pddlgym} bridges the gap between planning and learning by providing Gym-compatible environments generated from PDDL domain files, enabling evaluation of agents that combine symbolic and learned components. Multi-agent evaluation has also been supported by platforms such as PettingZoo~\citep{terry2021pettingzoo}, which provides a standardized interface for multi-agent reinforcement learning. These environments have been productive for advancing algorithms within their respective paradigms, but they do not support evaluation of how agents respond to changes that violate their learned or specified models---precisely the setting formalized in \cref{sec:open_world_environments}.

\paragraph{Generalization and domain transfer.}
A separate line of work addresses generalization across environments through domain randomization~\citep{tobin2017domain,peng2018sim} and procedurally generated environments. Benchmarks such as those developed by \citet{nichol2018gotta} evaluate transfer learning and few-shot learning in RL through procedurally generated levels that test whether agents generalize beyond their training configurations. The WordCraft environment~\citep{jiang2020wordcraft} incorporates natural language corpora to develop commonsense agents that can reason about novel situations. These approaches expose agents to a distribution of tasks during training, with the goal of improving robustness to variation at deployment. While effective for improving generalization within the training distribution, they do not explicitly evaluate how agents respond to post-training changes that introduce mismatch between the agent's models and the environment. The variability is incorporated into the training distribution rather than introduced as a structured, diagnosable deviation, and consequently these approaches do not exercise the adaptation capabilities that are central to open-world operation.

\paragraph{Out-of-distribution detection and novelty detection.}
Research on out-of-distribution (OOD) detection addresses the problem of identifying when an input differs from the training distribution. Methods such as confidence-based detection~\citep{hendrycks2017baseline}, open-set recognition~\citep{scheirer2013toward,bendale2016towards}, and autoencoder-based anomaly detection~\citep{abati2019autoregression} have been studied extensively in perception and classification settings, where novelty is defined in terms of input statistics. More recent approaches include vision transformer-based methods~\citep{lee2022anovit} and classifier-based novelty scoring~\citep{cheng2021learning}, which improve detection accuracy but remain focused on perceptual novelty. However, these methods focus on detection rather than adaptation: they identify that something has changed but do not provide mechanisms for the agent to modify its behavior in response. In the framework of \cref{sec:novelty}, detection addresses only one aspect of the agent's response to mismatch; the benchmarks introduced in this chapter must additionally support evaluation of how agents recover from and adapt to detected novelties.

\paragraph{Open-world novelty frameworks.}
The DARPA SAIL-ON program~\citep{senator2019science} catalyzed research on agents operating in environments with unanticipated novelty, motivating formal definitions of open-world novelty~\citep{boult2021towards,langley2020open} and agent architectures for novelty handling~\citep{muhammad2021novelty,goel2024neurosymbolic}. Within this broader effort, several evaluation environments have been developed. NovGrid~\citep{balloch2022novgrid} provides a flexible gridworld platform for studying novelty, built on MiniGrid, supporting configurable environment parameters and novelty injection. The Polycraft World AI Lab~\citep{goss2023polycraft} offers a richer, Minecraft-based environment for evaluating novelty detection and accommodation at scale, supporting multi-agent scenarios with cooperative and adversarial actors. NovelCraft~\citep{feeney2022novelcraft} provides a dataset and benchmark for visual novelty detection and discovery in open worlds.

Model-based approaches to novelty adaptation have also been evaluated in dedicated environments. \citet{klenk2020hydra} present HYDRA, a model-based novelty adaptation framework that uses a domain-independent planner and learns environment models that can be updated when novelties are introduced. \citet{stern2022hydra} extend this approach with a domain-independent planning framework evaluated in the Angry Birds domain. The Tonic library~\citep{pardo2020tonic} provides a configurable benchmarking framework for deep RL but is limited to closed-world settings and does not support novelty injection.

These environments represent important steps toward open-world evaluation, but they address different aspects of the problem than the benchmarks introduced in this chapter. NovGrid focuses on environment configurability but does not provide a formal taxonomy that connects novelty types to their impact on agent task performance, as developed in \cref{sec:novelty_taxonomy}. The Polycraft environment supports large-scale evaluation but is tightly coupled to a specific domain and agent interface, making it difficult to evaluate agents of different architectural paradigms under identical conditions. NovelCraft targets visual novelty detection rather than the full cycle of detection, adaptation, and recovery. None of these platforms provide a unified interface that seamlessly supports symbolic planning agents, reinforcement learning agents, and hybrid neurosymbolic architectures within the same environment.

\paragraph{Gap addressed by this chapter.}
The two benchmarks introduced in this chapter address gaps that cut across these lines of work. \NovelGridworlds{} and \NovelGym{} support controlled injection of diverse classes of mismatch, organized according to the taxonomy of novelty by domain element and task impact developed in \cref{sec:novelty}. They accommodate agents across multiple paradigms---symbolic planning, reinforcement learning, and hybrid architectures---through flexible observation interfaces, including LiDAR-like sensors, image-based local views, and symbolic PDDL representations. The \NovelGym{} ecosystem further supports multi-agent scenarios through its PettingZoo-based architecture~\citep{terry2021pettingzoo} and provides modular novelty injection through composable environment transformations. Both environments introduce evaluation metrics that capture not only task success but the dynamics of adaptation: how quickly performance degrades upon novelty injection, how efficiently the agent recovers, and whether beneficial novelties are exploited. These capabilities are essential for evaluating the claims developed in \cref{sec:abstractions}: that structured abstractions over actions and interactions determine whether agents can adapt efficiently and generalize effectively under novelty.
